% SOLO P2.
\section{Conclusiones}
\label{sec:conclusiones}

En este trabajo se evaluó la identificación del abandono de clientes en Telco Customer Churn mediante
un protocolo reproducible libre de fugas de información, comparando un clasificador de referencia con
regresión logística y árbol de decisión. Ambos modelos superaron ampliamente al \emph{baseline} en la
detección de abandonos (\emph{Recall} de 0.781 y 0.722 frente a 0.000), destacando la regresión
logística por su mejor ordenamiento global y menor brecha entre entrenamiento y validación, aunque el
conjunto de datos impone un techo informativo compartido en los casos de perfil estable.

Respecto a las preguntas de investigación planteadas, se concluye lo siguiente:

\begin{enumerate}
    \item \emph{Características con diferencias de abandono.} El contrato, la antigüedad, el tipo de
    internet, el método de pago y los servicios de seguridad y soporte muestran las mayores
    diferencias en el EDA. Ambos modelos confirman que \var{Contract}, \var{tenure} e
    \var{InternetService} concentran la mayor importancia por permutación.

    \item \emph{Contrato y antigüedad.} La diferencia por contrato persiste al controlar por
    antigüedad y demás variables: en la regresión logística, un contrato de dos años multiplica los
    \emph{odds} de abandono por 0.24 frente al mensual (y uno de un año, por 0.46), mientras que el
    árbol elige \var{Contract} como su primera división.

    \item \emph{Internet y servicios adicionales.} La fibra óptica se asocia con mayor abandono
    incluso dentro del contrato mensual (55.5\% frente a 28.1\% en el árbol; OR de DSL frente a
    fibra: 0.50), mientras que \var{OnlineSecurity} y \var{TechSupport} se asocian con menor riesgo
    (OR de 0.66 y 0.69).

    \item \emph{Calidad, representatividad y fuga de información.} Los 11 faltantes de
    \var{TotalCharges} se imputaron con cero, se excluyó \var{customerID}, los preprocesadores se
    ajustaron solo con entrenamiento y el conjunto de prueba se usó una sola vez. La
    representatividad sigue siendo la limitación principal por tratarse de una muestra ficticia y sin
    marcas temporales.

    \item \emph{Capacidad para identificar abandonos.} El clasificador de referencia no detecta
    ninguno de los 374 abandonos de prueba. La regresión logística detecta 292 (\emph{Recall} 0.781,
    F1 0.617, ROC-AUC 0.841) y el árbol, 270 (0.722, 0.599 y 0.818), con una \emph{Precision} cercana
    a 0.51. La exactitud casi no cambia respecto de la referencia (0.743 y 0.744 frente a 0.735), lo
    que confirma que no es la métrica adecuada para clases desbalanceadas.

    \item \emph{Desempeño e interpretabilidad.} La regresión logística supera al árbol en
    \emph{Recall} y ROC-AUC; en F1 la diferencia no es concluyente. Se interpreta mediante razones de
    \emph{odds} globales, mientras que el árbol elegido alcanza 132 hojas (legible solo en sus
    primeros niveles), aunque un árbol simple de 8 hojas logra un desempeño estadísticamente
    equivalente (F1 = 0.608 en validación).
\end{enumerate}

Como principales aprendizajes, el tratamiento del desbalance influyó mucho más que la fuerza de la
regularización $C$; cuantificar la incertidumbre mediante \emph{bootstrap} evitó presentar como
ventajas diferencias de F1 que están dentro del ruido muestral; restringir el tamaño mínimo de hoja
evitó la memorización del árbol (0.996 en entrenamiento frente a 0.496 en validación sin podar); y
fijar el protocolo antes de modelar garantizó resultados comparables en todo el equipo.

Como trabajo futuro se propone: (1) seleccionar hiperparámetros con validación cruzada estratificada
para reducir la dependencia de una partición única; (2) optimizar el umbral de decisión a partir de
costos reales de retención y pérdida; (3) calibrar las probabilidades estimadas mediante regresión
isotónica o escalamiento de Platt; (4) evaluar variables derivadas, interacciones y modelos de
ensamble como \emph{random forest} o \emph{gradient boosting}; y (5) validar con datos temporales de
un operador real monitoreando los errores por grupo sensible.
